\def\stat{l2-8}





\def\tit{АСПЕКТЫ ПРОСТРАНСТВЕННОЙ СОГЛАСОВАННОСТИ ГЕОГРАФИЧЕСКОЙ


ИНФОРМАЦИОННОЙ СИСТЕМЫ}








\def\titkol{Аспекты пространственной согласованности ГИС} 





\def\autkol{С.\,К. Дулин, И.\,Н. Розенберг, В.\,И. Уманский}





\def\aut{С.\,К. Дулин$^1$, И.\,Н. Розенберг$^2$, В.\,И. Уманский$^3$}





\titel{\tit}{\aut}{\autkol}{\titkol}





%{\renewcommand{\thefootnote}{\fnsymbol{footnote}}\footnotetext[1]


%{Работа выполнена при поддержке РФФИ, грант 10-01-00480.}}





\renewcommand{\thefootnote}{\arabic{footnote}}


\footnotetext[1]{Открытое акционерное общество <<Научно-исследовательский и 


проектно-конструкторский институт информатизации, автоматизации и связи на 


железнодорожном транспорте>>  (ОАО <<\mbox{НИИАС}>>); Институт проблем информатики Российской 


академии наук, s.dulin@gismps.ru}


\footnotetext[2]{Открытое акционерное общество <<Научно-исследовательский и 


про\-ект\-но-кон\-ст\-рук\-тор\-ский институт информатизации, автоматизации и связи на 


железнодорожном транспорте>>  (ОАО <<\mbox{НИИАС}>>}


\footnotetext[3]{Закрытое акционерное общество <<ИнтехГеоТранс>>}


     


\Abst{В географических информационных системах (ГИС), где управ\-ле\-ние качеством 


геоданных зависит от специфики пространственных объектов, проблема согласованности 


приобретает дополнительные сложности. 


      В~статье описаны особенности пространственной согласованности географических 


наборов данных в векторном формате, существенные при устранении ошибок отображения и 


классификации. Рассматривается три вида ошибок, которые влияют на соответствующие 


виды согласованности: структурные, гео\-мет\-ри\-че\-ские и топосемантические. Каждый вид 


ошибки требует разработки специальных процедур проверки и коррекции. Эти процедуры 


представлены  в работе в виде общей структуры.}


      


      \KW{геоинформационная система; топология геоданных; структуризация; 


согласованность ГИС}


      


      %Keywords: geoinformation system, topology of geodata, structurization, coordination 


%GIS.


      


                        \vskip 12pt plus 9pt minus 6pt








      \thispagestyle{headings}





      


            \label{st\stat}


            


            \vspace*{-9pt}


      


\section{Введение}





    Качество данных~--- ключевой вопрос любой информационной системы. 


В~ГИС геометрические 


особенности геоданных определяют формирование комплекса управления их 


качеством. Комиссией пространственного качества данных ICA~[1] были 


определены семь компонентов пространственного качества геоданных: их 


происхождение, позиционная точность, атрибутивная точность, завершенность, 


логичность~[2], семантическая точность~[3] и временн\'{а}я зависимость. 


    


    В данной работе анализируются логическая согласованность и 


семантические аспекты точности представления геоданных. 


    


    Независимо от происхождения (оцифровка карты, аэрофотосъемка или 


данные системы GPS~--- Global Positioning System) 


получающиеся географические наборы данных должны 


быть непротиворечивыми, чтобы использоваться в процессе пространственного 


анализа и гарантировать надежность решений. 


% 


    Однако во многих существующих наборах геоданных ощущается нехватка 


геометрического и топологического структурирования, что неизбежно 


приводит к ошибкам, не обеспечивая должным образом надежность 


результатов запросов, анализа или рассуждения. 


    


    Представленный ниже подход к улучшению пространственной 


со\-гла\-со\-ван\-ности существующих наборов геоданных в векторном формате 


позволяет улучшить как согласованность между объектами базы геоданных и 


реальными объектами (стрелка~1 на рис.~\ref{f1-2-8}), так и внутреннюю 


согласованность базы геоданных самой по себе (модели и структуры, стрелка~2 


на рис.~\ref{f1-2-8}). 


    


\begin{figure} %fig1


\vspace*{1pt}


\begin{center}


\mbox{%


\epsfxsize=81.44mm


\epsfbox{ipi-21(2)-8-1.eps}


}


\end{center}


\vspace*{-9pt}


\Caption{Проверка согласованности существующих ГИС


\label{f1-2-8}}


\end{figure}





В работе используется понятие топосемантической (topo~--- от 


топологического) согласованности, которая является частью логической 


согласованности, впервые определенной в~[2]. Топосемантическая 


согласованность касается правильности топологических отношений между 


двумя объектами согласно их семантике. Например, здание в другом здании~--- 


это, конечно, ошибка, тогда как здание на участке земли~--- не ошибка, хотя в 


обоих случаях используется отношение <<многоугольник в многоугольнике>>. 


Поэтому использование семантики, присущей каждому объекту, обязательно, 


чтобы обеспечить корректность каждого отношения. 





\section{Пространственный контекст согласованности}


    


    В настоящее время ГИС все более вовлечены в процессы, основанные на 


логическом рассуждении и принятии решений с учетом пространственных 


особенностей объектов. Картография больше не является основной зоной 


функциональности ГИС. 


    


    Существующие наборы геоданных содержат ошибки~[4] (и особенно 


ошибки, не видимые в рабочем масштабе), которые не нарушают 


визуализацию, но препятствуют или делают вообще невозможным проведение 


пространственного анализа. Более того, пространственным моделям ГИС 


недостает структуризации, чтобы обеспечить пространственный анализ. Таким 


образом, методология проверки и коррекции согласованности жизненно 


необходима для реализации пространственного анализа средствами ГИС. 


    


    Предлагаемое исследование имеет дело только с базами геоданных в 


векторном формате. Тогда для разработки метода проверки и коррекции 


согла\-со\-ван\-ности достаточно, чтобы были доступны только базы геоданных, 


модели геоданных и некоторые спецификации. 


    


    Проверка пространственной согласованности требует разработки процедур 


обнаружения несогласованности и определения ошибок согласованности. 


Ошибка может иметь последствия различного уровня. Некоторые ошибки 


могут отключить процесс анализа и тем самым быть своевременно 


обнаруженными, тогда как другие ошибки  могут привести к ошибочным 


конечным результатам. В~этом случае трудно выяснить, надежен результат 


или нет. 


    


    Получение геоданных системой может быть выполнено по-разному. 


Исследование различных методик позволяет выдвинуть предположения 


относительно происхождения ошибок, сопровождающих данный процесс. 


В~ГИС обычно используют следующие методы: 


    \begin{itemize}


\item \textit{сбор координатных геоданных непосредственно на экране} (с 


устройством позиционирования): в таком случае оценка только визуальна и 


субъективна (это вызывает существенную погрешность); инструментальные 


средства проверки, которые гарантируют правильный ввод, встречаются редко 


и, кроме того, они не всегда могут обеспечить полную проверку; 


\item \textit{сбор координатных геоданных, набираемых непосредственно с 


клавиатуры}: возникают проблемы законности геоданных (легитимность 


происхождения, неясная степень точности и качества), обнаруживаются 


ошибки ввода;


\item \textit{автоматическое пополнение}: возникают проблемы извлечения и 


дешифрирования геоданных при интерпретации аэрофотосъемки;


\item \textit{получение геоданных, поступающих из других систем}: степень 


надежности и точности геоданных не всегда известны, ошибки могут 


произойти и в процессе передачи. Могут возникнуть семантические ошибки 


(основанные на топологических отношениях) из-за различий между этими 


системами (различные структуры геоданных или системы координат);


\item \textit{композиция (например, объединение) с другими объектами, 


существующими в базе геоданных}: распространение возможных ошибок 


используемых объектов; 


\item \textit{преобразование} в цифровую форму: погрешность точности 


позиционирования, ошибки позиционирования, многократные точки (при 


расположении слишком близко друг к другу), объекты, зафиксированные 


дважды. 


\end{itemize}





    Тем самым очевидно, что процессы сбора геоданных потенциально могут 


быть источниками ошибок. Некоторые из таких ошибок могут быть 


обнаружены прямой экспертизой (незамкнутый многоугольник, точка вне 


диапазона), в то время как для обнаружения других видов ошибок необходимо 


принять во внимание семантику объектов (два многоугольника, чье перекрытие 


может быть ошибкой в зависимости от реальных объектов, которые 


они представляют). Следовательно, можно определить три вида ошибок в 


зависимости от того, какая часть объектов рассматривается.


    


    Выявление каждого вида ошибок соответствует уровню представления 


объектов, которые должны быть проанализированы, чтобы установить их 


согласованность. Существует три уровня представления: 


    \begin{enumerate}[(1)]


\item структурный уровень (структуры геоданных), 


\item геометрический уровень (концептуальная модель), 


\item топосемантический уровень (значения объектов согласно 


топологическим отношениям). 


\end{enumerate}





Соответственно имеют место три вида ошибок. 


    


    Структурные ошибки возникают из-за некорректности структур 


геоданных. Структуры должны позволять хранить данные согласно выбранной 


модели.\linebreak\vspace*{-12pt}


\begin{floatingfigure}{60mm} %fig2


\vspace*{1pt}


\begin{center}


\mbox{%


\epsfxsize=55.589mm


\epsfbox{ipi-21(2)-8-2.eps}


}


\end{center}


\vspace*{-9pt}


\Caption{<<Плохой>> многоугольник с внут\-рен\-ним контуром


\label{f2-2-8}}


\vspace*{6pt}


\end{floatingfigure}





\noindent


 Иногда специфические характеристики, поддерживаемые моделью, не 


поддерживаются структурами геоданных, что может привести к ошибкам. 


Например, некоторые объекты ГИС, состоящие из нескольких 


многоугольников, хранятся как многоугольники с внутренними контурами~[5] 


(рис.~\ref{f2-2-8}). Несогласованности могут происходить здесь на стадии 


реализации выполнения различных процедур. Структура геоданных, которая не 


отражает модель достаточно достоверно, может привести к структурным 


ошибкам. 


    





    


    Геометрические ошибки происходят из-за неправильной интерпретации 


геометрической части объектов (формы и позиции). Модель геоданных должна 


давать адекватное представление мира. Однако некоторые особенности 


реальной геометрии объектов не всегда хорошо фиксируются выбранной 


моделью. Например, многоугольник должен быть замкнут; следовательно, если 


он не\-замк\-ну\-тый~--- это геометрическая ошибка (если модель не определяет 


многоугольник как замкнутый объект, то тогда эта модель противоречива). 


    


    Топосемантические ошибки связаны со значениями характеристик 


реальных объектов, представленных в базе геоданных. Их следует 


обнаруживать и рассматривать на основе анализа пространственных 


отношений. 





\section{Проверка согласованности геообъектов}


    


    Географическое информационное моделирование сопряжено с 


определенными проблемами, связанными с пространственными признаками 


геообъектов. Реальные объекты, типа зданий или озер, характеризуются 


формой, местоположением, отношениями с другими объектами и семантикой. 


Форма и местоположение являются геометрическими признаками. В~процессе 


моделирования географической информации в первую очередь следует 


обратить определенное внимание именно на эти два вида признаков. Наряду с 


геометрическими признаками объектов очень важны пространственные 


отношения. Пространственные отношения обычно группируются в три 


категории~[5]: 


    \begin{enumerate}[(1)]


\item топологические отношения, которые являются инвариантными при 


топологических преобразованиях рассматриваемых объектов;


\item метрические отношения в терминах расстояний и направлений;


\item отношения пространственного порядка, определяющие порядок между 


объектами в зависимости от точки наблюдения. 


\end{enumerate}





    Проверка пространственной согласованности должна учитывать все 


особенности этих категорий. Цель проверки состоит в том, чтобы 


гарантировать, что геометрический объект и пространственные признаки 


геообъектов правильно обработаны базой геоданных. На рис.~\ref{f3-2-8} 


показана схема метода проверки пространственной согласованности. 








    Список свойств объектов фиксируется на основе наблюдений окружающей 


среды и исследования наиболее используемых пространственных моделей 


геоданных. Состав списка, указанного на рис.~\ref{f3-2-8}, будет рассмотрен в 


следующем разделе. Выбирая свойства из этого списка согласно модели и 


качественным спецификациям базы геоданных, можно сформировать список 


свойств, необходимых для каждой ГИС. Сформированный список свойств 


используется для проверки геометрической согласованности набора геоданных. 


    


    Топологические отношения между объектами могут быть описаны на 


основе модели 9~пересечений~[6, 7]. Топологические ограничения 


целостности, в свою очередь, определяются с помощью топологических 


отношений. Построенные таким образом ограничения можно затем 


использовать, чтобы проверить топосемантическую согласованность набора 


геоданных. 


    


    Данные в базе геоданных хранятся в виде определенных структур. Эти 


структуры соответствуют принятой модели, но иногда требуются специальные 


приемы программирования, чтобы обрабатывать специфические признаки 


(например, многоугольники с внутренними контурами (см.\ рис.~\ref{f2-2-8})). Этот 


вид ошибки связан с несоответствием модели и структуры геоданных и не 


может быть определен общепринятым способом. Однако такие проблемы 


должны быть зафиксированы и разрешены в рамках проверки структурной 


согласованности. 


\begin{figure} %fig3


\vspace*{1pt}


\begin{center}


\mbox{%


\epsfxsize=118.203mm


\epsfbox{ipi-21(2)-8-3.eps}


}


\end{center}


\vspace*{-9pt}


\Caption{Схема метода проверки согласованности


\label{f3-2-8}}


\end{figure}





\subsection{Геометрическая согласованность}


    


    Назначение модели геоданных состоит в том, чтобы  представить 


окружающую среду в базе геоданных. Это представление должно быть 


простым и отражать существенные для решения поставленных задач 


особенности реальных объектов. Геометрическое моделирование объектов 


должно удовлетворять двум главным требованиям: 


    \begin{enumerate}[(1)]


\item хорошее математическое представление объектов, 


\item хорошее описание объектов. 


\end{enumerate}





    В существующих ГИС модели отвечают этим требованиям при различных 


уровнях полноты представления геоданных. 


    


    Возникает первый вопрос процесса проверки согласованности: <<Модель 


геоданных отражает все релевантные особенности реальных объектов?>> 


    %


    Ответ, конечно: <<нет>>. 


    


    Прежде всего, реальный мир содержит 


слишком много отображений, которые не могут быть зафиксированы в модели 


геоданных. Во-вторых, существующие модели были спроектированы, чтобы 


удовлетворять информационным нуждам, спектр которых постоянно 


расширяется. 


    


    Следовательно, модели геоданных должны быть улучшены добавлением 


свойств, которые фиксируют пространственные характеристики реальных 


объектов. 





\begin{figure} %fig4


\vspace*{1pt}


\begin{center}


\mbox{%


\epsfxsize=126.623mm


\epsfbox{ipi-21(2)-8-4.eps}


}


\end{center}


\vspace*{-9pt}


\Caption{Свойства геометрических объектов и связей между ними


\label{f4-2-8}}


\end{figure}


    


    Свойство~--- это правило о форме или о позиции нужного объекта, 


которое должно сохраняться. Каждый раз, когда правило не выполняется, 


обнаруживается геометрическая ошибка. На рис.~\ref{f4-2-8} показан ряд 


свойств, которые могут быть зафиксированы в базе геоданных. Каждое 


свойство присоединено к форме или к связи между двумя формами.





\smallskip





{\bfseries\textit{Свойства объектов}}





\smallskip





 p1~--- Точечная согласованность (диапазон значения на каждой оси). 





     p2~--- Точка или уникальность узла. 





     p3~--- Линия или уникальность кромки (края). 





     p4~---  Многоугольник или уникальность переднего края. 





 p5~---  Существование на уровне ссылок (это не пространственное свойство, 


поэтому оно не показано на рис.~\ref{f4-2-8}). 





p6.1~--- Незамкнутый объект.





 p6.2~--- Замкнутый объект. 





     p7~---  Несамопересекающийся (по границе). 





     p8~---  Связность. 


    


 p9~---  Значение края. 


    


 p10~---  Ориентация. 





 p10.1~--- Граничная ориентация. 





 p11~--- Пространственное покрытие рядом объектов. 





 p12~--- Неперекрывающиеся многоугольники или передние планы. 














\smallskip


{\bfseries\textit{Свойства связей между объектами }}





\smallskip





     pI~---  Две различные точки или узлы. 





     pII~---  Принадлежит, по крайней мере, двум объектам. 





 pIII~--- Принадлежит точно одному. 





 pIV~--- Все объекты набора различны. 





     pV~---  Объекты набора отсортированы. 





 pVI~---  Принадлежит точно двум объектам. 





 pVII~---    Принадлежит, самое большее, двум объектам. 





    В зависимости от модели геоданных можно построить соответствующий 


список свойств, подходящих для ГИС, выбирая свойства из списка. Этот 


процесс показан на рис.~\ref{f3-2-8} стрелкой, названной выбором.


    


    Модели, основанные на многоуровневом представлении, планарных 


графах и многоугольных моделях геоданных~--- самые популярные 


геометрические модели, используемые в ГИС (описание этих моделей можно 


найти в~[4, 8]). Они основаны на геометрических объектах, таких как 


точки (или узлы), линии (или дуги), многоугольники (или передние планы) и 


отношениях между такими объектами (например, линия~--- ряд точек, которые 


могут быть отсортированы или нет). 


    


    Структуры геоданных, используемые для реализации таких моделей, могут 


содержать некоторые свойства объектов в своем определении, т.\,е.\ объекты, 


внесенные в структуру, автоматически обладают ими. Следовательно, свойства, 


зафиксированные моделями геоданных, можно не проверять~[9]. 


    


    В представленной работе приняты во внимание только топологические 


ограничения. Топологическое отношение между двумя сущностями основано 


на анализе общих частей их форм. Достоверность любого топологического 


отношения основана на семантике обоих сущностей. Именно поэтому ошибки, 


исходящие из недостоверных топологических отношений между объектами, 


называют топосемантическими. 


    


    Существует несколько моделей для обработки топологических 


отношений~[6, 7]. Ниже представлены модель 9~пересечений и метод 


проектирования топологических ограничений целостности, основанный на этой 


модели. 


    


    Топологические ограничения целостности основаны, как уже говорилось 


выше, на топологическом отношении. Такие отношения описывают 


относительную позицию объектов в пространстве. 


    


    Топологическая модель 9 пересечений была достаточно хорошо 


изучена~[8, 10] и специфицирована М.\,Дж.~Эдженхофером~[5--7]. В~этой 


модели бинарные топологические отношения между двумя объектами~$A$ 


и~$B$ определены в терминах 9~пересечений  границы~$A$ ($\partial A$), 


внутренней части области~$A$ ($A^\circ$) и внешней части области $A$ ($A^-


$) с границей~$B$ ($\partial B$), внутренней частью области~$B$ ($B^\circ$) и 


внешней частью области~$B$ ($B^-$). Каждый объект~$A$ и~$B$ может быть 


точкой, линией или многоугольником. Определения каждой части каждого вида 


геометрического объекта следующие: 


    


$P$~--- точка: $P=\partial P =P^\circ$;





$L$~--- линия: $\partial L = \mbox{две\ концевые\ точки}\ L$;





$L^\circ = L - \partial L$;





$R$~--- многоугольник: $\partial R$\;=\;пересечение замыкания $R$ и 


замыкания внешней части~$R$;





$S^\circ$\;=\;объединение всех открытых множеств в~$R$.


    


    Для каждого пересечения вычисляется пустое ($\varnothing$) или непустое 


($\neg \varnothing$) значение и сохраняется в  матрице: 


    $$


    \begin{pmatrix}


\    \partial A\cap\partial B \ \, &\ \, \partial A\cap B^\circ \ \,&\ \,\partial A\cap B^-\ \\[6pt]


 \    A^\circ\cap\partial B \ \, &\ \, A^\circ\cap B^\circ \ \,&\ \, A^\circ\cap B^- \ \\[6pt]


  \   A^-\cap\partial B \ \,&\ \, A^-\cap B^\circ \ \,&\ \, A^-\cap B^-\ 


    \end{pmatrix}\,.


    $$





    Топологические ограничения целостности определены на основе 


топологических отношений, описанных моделью 9~пересечений. При этом 


топологическое отношение между двумя объектами~--- основная часть 


ограничения. Рассматривая форму объектов, можно вычислить все допустимые 


топологические отношения между двумя объектами (согласно модели 


9~пересечений). Рассматривая тип объекта, можно определить, какое 


топологическое отношение является непротиворечивым, а какое~--- 


противоречивым. 


    


    Топологическое ограничение может быть определено как 


\textbf{ОГРАНИЧЕНИЕ}\;=\;\textbf{(Класс\_сущностей~1, Отношение, 


Класс\_сущностей~2, Спецификация)}~--- совокупность двух географических 


объектов, топологического отношения между ними и спецификации, которая 


может быть одной из следующих:


    \begin{itemize}


\item запрещенное; 


\item по крайней мере,  $n$  раз; 


\item не больше, чем  $n$  раз; 


\item точно  $n$  раз. 


    \end{itemize}


    


    Спецификация <<Запрещенное>>~--- самая интересная и полезная для 


использования. Топологические ограничения целостности, определенные с 


использованием этой спецификации, нужны для конечных пользователей, 


чтобы описать топологические ситуации, которые они не хотят видеть в своей 


базе геоданных. 


    


    Модель~9 пересечений может быть применена ко всем видам 


геометрических объектов. 


    


    Рассмотрение в качестве основных сущностей точки, линии и 


многоугольника позволяет сформировать 6~групп возможных отношений: 


точ\-ка/точ\-ка, точ\-ка/ли\-ния, точ\-ка/мно\-го\-уголь\-ник, ли\-ния/ли\-ния, 


ли\-ния/мно\-го\-уголь\-ник, многоугольник/многоугольник. 


    


    Рассматриваемая модель позволяет определять только простые 


ограничения (отношение между двумя объектами). В~[11] определена модель для 


топологических ограничений целостности, также основанных на модели 


Эдженхофера. Предложенная там модель позволяет определять тот же самый 


вид ограничений и характеризовать сложные пространственные отношения, 


комбинирующие несколько отдельных ограничений. Но такая модель меньше 


устраивает конечных пользователей, так как такой способ определения 


ограничений требует определенных навыков составления логических условий. 


    


    Ограничения топологической целостности определяют правила, 


основанные на семантике сущностей базы геоданных. Как и список свойств, 


эти ограничения должны быть адаптированы к обрабатываемому набору 


геоданных. 


    


    Топологические ограничения определяются следующим списком 


операций: 


    \begin{itemize} %[1)]


\item выбрать первый класс объектов; 


\item выбрать второй класс объектов; 


\item выбрать отношение из предложенного 


списка; 


\item определить спецификацию. 


\end{itemize}


     


Примеры ограничений топологической целостности: 





ОГРАНИЧЕНИЕ1 (Дорога, Пересечение, Здание, Запрещенное); 





ОГРАНИЧЕНИЕ2 (Водовод, Соединение, Водопровод, Точно 2~раза). 





     Модель~9 пересечений приводит к 81~отношению между точками, 


линиями и многоугольниками~[6, 7]. Это число является реальным 


препятствием конструирования ограничений по двум причинам: 


     \begin{enumerate}[(1)]


\item слишком много отношений, чтобы именовать каждое из них; 


\item чтобы исключить конкретную ситуацию, нужно создавать несколько 


ограничений. 


\end{enumerate}





    По этим причинам топологические отношения целесообразно 


сгруппировать в подмножества~[8]. Использование подмножества в 


определении ограничений делает эту модель более практичной для 


пользователя. 


    


    Чтобы вычислять топологические отношения, необходимо вычислить 


матрицу 9~пересечений (или, по крайней мере, некоторые ее элементы), т.\,е.\ 


пересечения границ, внутренних и внешних областей двух объектов. Каждое 


ограничение затем может быть переведено в конъюнкцию проверок согласно 


отношению (или подмножеству отношения), включенному в него. 


    


    Каждое подмножество или каждое отдельное топологическое отношение 


может быть описано частичным отношением. Например, линия, пересекающая 


линию, определяется единственным элементом матрицы  9~пересечений: 


%  


    $A^\circ\cap B^\circ\hm=0$  и приводит только к одной проверке. Каждая 


проверка описывается предикатом~[8], основанным на размерности 


пересечения ($\varnothing$, 0, 1 или 2~--- первый столбец табл.~\ref{t1-2-8}). 


В~приведенном примере: ПЕРЕСЕЧЕНИЕ\_РАЗМЕРНОСТИ\_0 ($A^\circ, 


B^\circ$). Чтобы оценивать такие предикаты, требуются функции для 


на\-хож\-де\-ния границы, внутренней и внешней области объектов. Таким образом, 


более точное выражение для приведенного примера выглядит так: 


ПЕРЕСЕЧЕНИЕ\_РАЗМЕРНОСТИ\_0 (внутренняя область ($A$), внутренняя 


область ($B$)). 


 % 


    Например, одномерное пересечение многоугольника и линии является 


линией. 





\begin{table}\small


\begin{center}


\Caption{Пример топологических отношений


\label{t1-2-8}}


\vspace*{2ex}





\begin{tabular}{|l|c|c|}


\hline


\multicolumn{1}{|c|}{Тип предиката}& Объект 1 & Объект 2 \\


\hline


ПЕРЕСЕЧЕНИЕ\_РАЗМЕРНОСТИ\_2 &Многоугольник &Многоугольник\\ 


\hline


\multicolumn{1}{|l|}{\raisebox{-6pt}[0pt][0pt]{ПЕРЕСЕЧЕНИЕ\_РАЗМЕРНОСТИ\_1}} &Многоугольник & Линия \\


& Линия & Линия \\


\hline


\multicolumn{1}{|l|}{\raisebox{-18pt}[0pt][0pt]{ПЕРЕСЕЧЕНИЕ\_РАЗМЕРНОСТИ\_0}} &Многоугольник& Точка \\


&Линия & Линия \\


&Линия & Точка \\


&Точка &Точка \\


\hline


\multicolumn{1}{|l|}{\raisebox{-28pt}[0pt][0pt]{ПУСТОЕ\_ПЕРЕСЕЧЕНИЕ}}&Многоугольник &Многоугольник\\


&Многоугольник &Линия\\


&Многоугольник&Точка \\


& Линия &линия \\


& Линия & Точка \\


&Многоугольник& Точка \\


     \hline


     \end{tabular}


     \end{center}


     \vspace*{-6pt}


     \end{table}


    


    Рассмотрение четырех типов предикатов и трех видов геометрических 


объектов, 9~функций и 13~предикатов достаточно, чтобы проверить все 


возможные ограничения топологической целостности (см.\ табл.~\ref{t1-2-8}). 


Имеются в виду сле\-ду\-ющие 9~функций: граница (точка); внутренняя область 


(точка); внешняя область (точка); граница (линия); внутренняя область (линия); 


внешняя область (линия); граница (многоугольник); внутренняя область 


(многоугольник); внешняя область (многоугольник). 


    


    Если можно вычислить границу, внутреннюю и внешнюю области каждого 


вида геометрического объекта, то можно и проверить ограничения 


топологической целостности независимо от используемой модели геоданных. 





    \begin{figure} %fig5


\vspace*{1pt}


\begin{center}


\mbox{%


\epsfxsize=127mm


\epsfbox{ipi-21(2)-8-5.eps}


}


\end{center}


\vspace*{-9pt}


\Caption{Структура проверки и коррекции согласованности ГИС


\label{f5-2-8}}


\end{figure}





\section{Структура для обнаружения ошибок и коррекции}


    


    Структура проверки и коррекции пространственной согласованности на 


рис.~\ref{f5-2-8} представляет связи между различными частями 


полнофункциональной ГИС, чтобы можно было проверять наборы геоданных и 


исправлять пространственные ошибки. Будем рассматривать две части: база 


геоданных, которая содержит всю информацию относительно описания 


реального мира, и функциональные возможности ГИС, которая содержат ряд 


инструментальных средств доступа к данным и обработки наборов геоданных. 


     


     


    База геоданных имеет концептуальный и структурный уровни. 


Концептуальный уровень~--- единственный, который полностью находится в 


ведении конечного пользователя. Через выбранную модель этот уровень 


определяет логическую организацию  базы геоданных. 


    


    Структурный уровень описывает все структуры хранения геообъектов. 





\begin{table}[b]\small


\begin{center}


\parbox{262pt}{\Caption{Набор топологических ограничений целостности для ГИС городского хозяйства


\label{t2-2-8}}





}





\vspace*{2ex}





    \begin{tabular}{|l|c|c|c|}


\hline


\multicolumn{1}{|c|}{Сущность 1}&Отношение & Сущность 2 


&Спецификация\\


\hline 


Район& Наложение & Район & Запрещенное \\


& ВНУТРИ && Запрещенное \\


\hline


Квартал &ВНУТРИ & Район & Точно $n$ раз \\


\hline


Квартал &Наложение& Квартал & Запрещенное\\ 


& ВНУТРИ &&Запрещенное \\


\hline


Участок& ВНУТРИ & Квартал & Точно $n$ раз \\


\hline


Участок & Наложение & Участок &Запрещенное \\


&ВНУТРИ &&Запрещенное \\


\hline


Здание &ВНУТРИ & Участок & Точно $n$ раз \\


\hline


Здание &Наложение & Здание & Запрещенное\\ 


&ВНУТРИ && Запрещенное \\


\hline


\end{tabular}


\end{center}


\end{table}


    


    Географическая информационная система 


    содержит инструментальные средства манипулирования базой 


геоданных. Целесообразно различать два вида инструментальных средств: 


инструментальные средства обработки, т.\,е.\ процессы, которые вычисляют 


результаты запросов, анализа, вывода, и инструментальные средства 


интерфейса, которые обеспечивают конечного пользователя средствами 


формирования этих запросов. В~нашем случае этот интерфейс содержит 


интерфейс определения топологических ограничений, интерфейс коррекции 


топологических ошибок, который предлагает все возможные исправления для 


каждой ошибки, и интерфейс коррекции геометрических ошибок. Уровень 


средств обработки содержит методы доступа к данным (которые зависят от 


структур геоданных), ведение журнала ошибок и процесс проверки 


топологических ограничений (трансляция ограничений и проверка). 


    


    Первый шаг процесса обнаружения ошибок должен адаптировать 


структуру к рассматриваемому набору геоданных. При этом определенной 


адаптации подвергаются как семантический, так и геометрический уровни. 


    


    Первая стадия предполагает выбор из списка на рис.~\ref{f4-2-8} свойств, 


которые релевантны обрабатываемой базе геоданных. 


    


    Затем следует обратиться к геометрической модели геоданных, 


используемой в ГИС (свойства, очевидно, различным образом зависят от 


выбранной модели: многоуровневая, планарный граф). Мы можем далее 


получить информацию из существующего документа о качестве (спецификации 


качества), в котором есть указания о том, что нужно проверять, а что нет. 


    


    Полученный таким образом исчерпывающий список геометрических 


свойств (относительно модели и спецификаций) должен быть зафиксирован в 


базе геоданных. 


    


    В табл.~\ref{t2-2-8} представлен набор топологических ограничений 


целостности для ГИС городского хозяйства~[8, 12] для проверки 


согласованности кадастрового уровня ГИС (который является подмножеством 


всей базы геоданных). Основные объекты, находящиеся на этом уровне,~--- 


здания, участки, кварталы и районы. 


     





    


\subsection{Обнаружение ошибки}


    


    Каждый вид ошибки обрабатывается определенным способом. Проверка 


согласованности на предмет структурных ошибок осуществляется 


процедурами, предназначенными для конкретной базы геоданных. 


    


    Геометрические ошибки являются, по сути, более общими, так как они 


связаны со свойствами геометрических объектов, обычно используемых в 


моделях геоданных. Поэтому процессы проверки согласованности, 


разработанные для одной базы геоданных, могут быть многократно 


использованы в другой базе. Тем не менее, некоторые свойства затрагивают 


конкретные структуры геоданных, и процесс проверки, связанный с ними, 


должен быть адаптирован к этим структурам. 


    


    Топосемантические ошибки зависят от общих понятий. При условии, что 


определенные функции доступа к данным были реализованы для базы 


геоданных, представленные процессы проверки могут быть применены к 


любой ГИС. 


    


    Вообще говоря, процессы проверки сталкиваются с двумя трудностями: 


проблема полноты всех видов возможных ошибок и проблема доказательства 


завершенности каждого процесса проверки: Действительно ли обнаружены все  


объекты, содержащие эту ошибку? Можно ли считать, что действительно 


обнаружены  некоторые хорошие объекты? 


    


    Структурные ошибки не могут быть установлены априори. Они возникают 


из-за некорректности структур геоданных, используемых в базе геоданных, 


поэтому список этих ошибок не может быть определен в начале обработки. 


    


    Такие ошибки не всегда могут выявляться в результате проверки свойств. 


Они часто связаны с использованием приемов программирования для 


обработки случаев, которые не определены моделью геоданных (например, для 


многоугольников с внутренними контурами, которые обрабатываются как 


сложные многоугольники). Набор приемов, обрабатывающих геометрические 


представления, компенсирует слабость модели. Напротив, более полная 


модель, с сильными, конструктивными и ясно определенными свойствами 


(такая как планарный граф) сама сокращает эти структурные ошибки. 


    


    Определение общих процессов проверки и коррекции этого вида ошибки 


выходит за пределы темы этой работы, но является предпосылкой общего 


процесса улучшения согласованности. Поэтому ниже представлен только один  


пример типичной структурной ошибки и способ ее обработки. 


    


    Самый интересный случай структурной ошибки касается многоугольников 


с внутренними контурами. Эта ошибка является следствием того, что простые и 


сложные многоугольники обычно не различаются в общей структуре 


геоданных. Отверстия описываются как часть границы объекта и связываются с 


внешними линиями двумя вспомогательными сегментами или 


вспомогательными линиями (с противоположными направлениями, чтобы 


гарантировать замыкание многоугольника (рис.~\ref{f6-2-8})): их называют 


соединяющими сегментами. 





\begin{figure} %fig6


\vspace*{1pt}


\begin{center}


\mbox{%


\epsfxsize=62mm


\epsfbox{ipi-21(2)-8-6.eps}


}


\end{center}


\vspace*{-9pt}


\Caption{Многоугольник с внутренним контуром


\label{f6-2-8}}


\end{figure}


    


    Таким образом, возникает два вида проблем, проистекающих из этого 


специфического средства представления: 


    \begin{enumerate}[(1)]


\item неточность нанесения этих соединяющих сегментов и погрешности их 


координат могут привести к ошибочным многоугольникам 


(самопересекающиеся границы~--- рис.~\ref{f7-2-8},\,\textit{a}, или отверстие 


больше не распознается как отверстие~--- рис.~\ref{f7-2-8},\,\textit{б}). 


    


    В первом случае (см.\ рис.~\ref{f7-2-8},\,\textit{a}) внутренняя и наружная 


границы имеют то же самое направление вращения (эта ошибка приводит к 


неправильному результату, когда вычисляется, например, поверхность объекта, 


\begin{figure} %fig7


\vspace*{1pt}


\begin{center}


\mbox{%


\epsfxsize=114.154mm


\epsfbox{ipi-21(2)-8-7.eps}


}


\end{center}


\vspace*{-9pt}


\Caption{Ошибки соединения сегментов


\label{f7-2-8}}


\end{figure}


тогда как во втором случае (см.\ рис.~\ref{f7-2-8},\,\textit{б}) никакое свойство в 


принципе не нарушено (тем не менее, это~--- действительно ошибка); 


\item вторая причина возникновения ошибок, сгенерированных этими 


объектами, непосредственно связана с различными назначениями этих 


соединяющих сегментов. Фактически они не принимают участие в описании 


объекта, в то время как на практике многочисленные многоугольники 


представлены с невидимыми вспомогательными сегментами, которые 


действительно описывают объекты, как показано на рис.~\ref{f8-2-8} (это 


обусловлено картографическими инструментальными средствами 


программного обеспечения, задействованными для того, чтобы избежать 


прорисовывания одной и той же линии несколько раз). Очень важно делать 


различие между этими двумя видами вспомогательных сегментов так, чтобы 


обрабатывать различные соединяющие сегменты разными способами, а не 


принимать во внимание их во всех случаях одинаково, выполняя 


геометрические операции (исчисление периметра, буферной зоны и~др.), 


которые могут привести к ошибочным результатам. 


    \end{enumerate}


    


    Геометрический вид ошибок может быть охарактеризован как неточность 


соотнесения объекта со свойствами принятой модели. 


    


    Выбор процедуры проверки зависит от модели геоданных. Для 


большинства свойств может быть определен специальный алгоритм, 


пригодный почти для всех сущностей модели (например, чтобы гарантировать 


замкнутое выражение многоугольника, достаточно иметь три алгоритма в 


зависимости от того, как определена его граница, используя точки, сегменты 


или дуги). Такие проверки~--- вопрос вычислительной геометрии. Без 


классификации векторных моделей геоданных в ГИС невозможно 


сформировать полный список алгоритмов проверки, однако во многих простых 


случаях может быть произведена обработка геоданных известными 


алгоритмами~[9].


    


    Еще один распространенный вид ошибки связан с искажением 


многоугольников и возникновением бесполезных точек. Этот вид ошибки 


возникает в процессе оцифровки. Из-за проблем с устройством 


позиционирования вмес\-то единственной точки может быть нанесено несколько 


точек, что приводит к появлению многоугольника с неправильными углами. 


Ошибка в этом случае может быть обнаружена либо благодаря условию 


уникальности точки на линии или в многоугольнике (свойство~pIV с допуском 


на минимальное расстояние между двумя точками), либо она может быть 


обнаружена с помощью условия несамопересекаемости многоугольника 


(свойство~p7). 


    


    Трудность в этом случае состоит в том, чтобы различать точки, которые 


определяют реальную кривую, и действительно бесполезные <<шумовые>> 


точки.


    


    И форму, и семантику объектов следует проверять на наличие таких 


ошибок. Форма позволяет определять топологическую сцену, а семантика 


позволяет устанавливать достоверность сцены (посредством ограничений 


топологической целостности). Проверка на наличие топосемантических 


ошибок требует вы\-чис\-ле\-ния топологического отношения. Чтобы упростить 


проблему, обычно рас\-смат\-ри\-ва\-ют\-ся только точки, линии и многоугольники 


(или области). 


    


    Например, пусть определено следующее, характерное для большинства 


ГИС топологическое ограничение целостности: 


    $\mathrm{(Река}$, $\mathrm{Пересечение}$, $\mathrm{Река}$, $\mathrm{Запрещенное})$.


Оно будет оттранслировано в: \textit{если} 


ПЕРЕСЕЧЕНИЕ\_РАЗМЕРНОСТИ\_0 (внутренняя область (Рекаl), внутренняя 


область (Река2)), \textit{тогда обнаружено} <<Несогласованность>>. 


    


    Каждый раз, когда обнаружена несогласованность, сцена сохраняется в 


журнале ошибок, который будет использоваться в процессе исправления 


геоданных. 


    


    Алгоритмы проверки ошибок обнаруживают несколько видов ошибок, 


позволяющих оценить число объектов, которые нужно исправить. 





\subsection{Исправление ошибок}


    


    В общем случае существует немного возможностей для автоматического 


исправления ошибок, потому что эти исправления могут быть сделаны, только 


если процесс обнаружения ошибок завершен и если существует только одно 


очевидно возможное исправление. В~других случаях необходимо исправление 


вручную. Это связано с тем, что конечному пользователю следует проверить, 


что ошибка~--- не исключение или что исправление нуждается в 


интерпретации семантики объекта, которую компьютер не может сделать. Тем 


самым целесообразно обеспечить пользователя интерфейсом, с помощью 


которого будут предлагаться различные сценарии исправлений и 


предотвращаться появление новых ошибок (полуавтоматическое исправление). 


    


    Что касается структурных ошибок, то процессы исправления будут сильно 


зависеть от рассматриваемой базы геоданных. Для случаев геометрических и 


семантических ошибок некоторые общие идеи, конечно, могут быть применены 


к любым ГИС. Однако способ обработки исправлений будет обязательно 


зависеть от структур геоданных. 


    


    Структурные ошибки непосредственно связаны с выбором реализации. 


Таким образом, они должны быть исправлены по отношению к структурам 


геоданных, и методы исправления должны быть определены для каждого 


случая структурных ошибок. 


    


    Что касается случая многоугольников с внутренними контурами, то для 


них имеет смысл определить новый стиль линии, чтобы отличать соединяющие 


сегменты от вспомогательных (так как они имеют различные назначения), и 


объединять концевые точки таких сегментов, когда это необходимо. 


    


    Большинство геометрических ошибок может быть сведено к проблеме, 


касающейся точек объектов. Поэтому должны быть ясно определены 


различные возможности обработки точек. Следующий список представляет 


основные операции, которые могут быть применены к точкам: 


    \begin{itemize}


\item добавление новой точки; 


\item удаление точки; 


\item слияние двух точек; 


\item проектирование точки на сегмент; 


\item изменение координат существующей точки. 


\end{itemize}





    Исправление многоугольника с неправильными углами (с бесполезными 


точками) может быть разделено на три части, как показано на рис.~\ref{f8-2-8}: 


    \begin{enumerate}[(1)]


\item вычисление лучшего местоположения вершины; 


\item проектирование; 


\item удаление бесполезных точек. 


\end{enumerate}





\begin{figure} %fig8


\vspace*{1pt}


\begin{center}


\mbox{%


\epsfxsize=120.198mm


\epsfbox{ipi-21(2)-8-8.eps}


}


\end{center}


\vspace*{-9pt}


\Caption{Исправление угла неправильного многоугольника


\label{f8-2-8}}


\end{figure}





    Рисунок~\ref{f8-2-8} иллюстрирует только самый простой процесс, 


который может быть применен к этому виду ошибки. Можно использовать 


несколько вариантов, чтобы найти лучшее местоположение для вершины, как-то: 


центр тяжести точек, пересечение двух значимых сегментов, использование 


других объектов. 


    


    Только ограничения топологической целостности, использующие 


спецификацию <<\textit{запрещенное}>>, могут привести к 


полуавтоматическим исправлениям. В~этом случае ошибка определена как 


запрещенное топологическое отношение между двумя объектами, поэтому 


способ исправления ошибки будет состоять в том, чтобы изменить 


топологическое отношение между этими объектами. Будет рассчитан ряд 


сценариев исправления с использованием нескольких видов изменений к обоим 


объектам, входящим в запрещенное топологическое отношение (вместе или 


один за другим). 





Предлагаются следующие изменения: 


    \begin{itemize}


\item модификация объектов: 


\begin{itemize}


\item перемещение объектов; 


\item изменение формы объектов; 


\end{itemize}


\item удаление одного объекта; 


\item разбиение объекта (создание нового объекта). 


\end{itemize}


    


    Расчет и предлагаемые исправления сценариев имеют два основных 


преимущества. Первое~--- это облегчение и ускорение работы конечного 


пользователя. Второе~--- управление процессом исправления гарантирует, что  


исправление не создаст новую ошибку. 


    


    \textit{Перемещение объекта} гарантирует, что поверхностная область 


обоих объектов остается неизменной. Один из двух объектов, входящих в 


запрещенное отношение, перемещается согласно основному направлению, пока 


топологическое отношение не изменится. 


    


    \textit{Изменение} формы означает перемещение части объекта, оставляя 


при этом другую часть неизменной. Цель такого исправления состоит в том, 


чтобы изменить топологическое отношение между двумя объектами, не 


меняя отношения с другими объектами набора геоданных. Корректировки 


будут сделаны некоторым соответствующим алгоритмом, который привяжет 


характерные точки одного объекта к характерным точкам другого объекта. 


    


    \textit{Удаление одного объекта} полезно, когда объект был оцифрован 


дважды. Тогда могут быть найдены два объекта, очень близко расположенные 


друг к другу. Однако удаление объекта не является простым мероприятием в 


базе геоданных из-за всех существующих ссылок на удаленные объекты 


(например, пространственная индексация). В~этом случае удаление~$A$ и 


сохранение~$B$ означает замену~$A$ на~$B$ и может быть реализовано 


проставлением флажка, означающего, что нужно всегда использовать 


объект~$B$ вместо~$A$ при любой ссылке на~$A$. Полное удаление объектов 


(и обновление всех ссылок) может быть выполнено впоследствии пакетной 


обработкой. 


    


    \textit{Разбиение одного объекта} на два новых объекта позволяет 


поддерживать планарность карты. Единственное условие проверки состоит в 


том, что два новых топологических отношения отличаются от предыдущего. 


Такая коррекция может быть предложена, когда запрещенное отношение 


возникло из-за того, что один из двух объектов покрывает частью своей 


внутренней области внутреннюю область или границу другого объекта. 


\pagebreak


    


    Тогда коррекции могут быть следующие: 


    \begin{itemize}


\item разбить один из двух объектов на несколько частей;


\item создать новый объект, основанный на общей части объектов, и удалить его 


от каждого из объектов. 


\end{itemize}





    Для каждой топологической ошибки создается список возможных 


коррекций. Последний шаг процесса коррекции должен выбрать приемлемую 


коррекцию. 


    


    Довольно часто используемые структуры геоданных основаны на 


многоуровневой модели. В~этом случае самые сложные проблемы 


заключаются в ис\-прав\-ле\-нии семантических ошибок, поскольку они основаны 


на топологии объектов. В~худшем случае конечному пользователю придется 


использовать собственные инструментальные средства фиксации и 


модификации, обеспечиваемые программно. Однако следует помнить, что 


такой метод может привести к появлению новых ошибок в базе геоданных. 


    


    Расширение описанной выше проблемы касается геоданных, которые 


могли быть обновлены извне (кадастровые данные, например). В~этом случае, 


поскольку модификации не были сделаны в системе, они не могли быть 


скопированы на другие уровни геоданных системы. Кроме того, журнал 


коррекций не мог быть доступен. Проблема состоит в том, чтобы заменить 


данные других уровней в правильном месте в соответствии с теми данными, 


которые уже были изменены. Таким образом, для решения этой проблемы 


необходимо предложить инструментальные средства для автоматического 


выполнения этих операций.


    


    Процессы проверки часто используют допуск, который в общем случае 


является максимальным расстоянием между точками, но может иметь и другой 


геометрический формат (угол между двумя сегментами)~[10]. Любой 


алгоритм, который использует допуск, учитывает выбранные ограничения. 


Фактически, допуск должен быть определен как абсолютное значение: ниже 


этого значения объект корректен (относительно проверяемого свойства), при 


превышении этого значения он некорректен. Такое четко определенное 


значение практически не может быть найдено; поэтому процесс обнаружения, 


который использует жесткий допуск, никогда не будет завершен: если значение 


допуска является слишком большим, ошибки могут остаться 


необнаруженными, и напротив: если это значение является слишком 


маленьким, могут быть пропущены объекты, которые не ошибочны. Таким 


образом, допуск должен устанавливаться соответственно проверяемым 


объектам и с точностью до набора геоданных, так, чтобы процесс обнаружения 


был завершенным, насколько это возможно. 


    


    Цель любой коррекции состоит в том, чтобы сохранить достигнутое 


качество в течение всего процесса коррекции и после его завершения. 


    


    При этом, прежде всего, различные процессы коррекции не должны 


привести к возникновению новых ошибок любого вида, т.\,е.\ применяемые 


алгоритмы должны гарантировать корректность исправленных геоданных 


относительно всех других видов ошибок. Таким образом, эти процессы должны 


выполняться только единожды. 


    


    Следует также обратить внимание на порядок коррекций. Поскольку 


некоторые ошибки могут быть исправлены только после того как будет решен 


ряд других проблем, процессы коррекции, конечно, должны упорядочиваться: 


тогда мы можем предположить, что будем исправлять ошибки, касающиеся 


только одного объекта (структурные и геометрические ошибки) до исправления 


тех ошибок, которые касаются нескольких объектов (семантические ошибки). 


Обратный метод привел бы к путанице с частями, которые не соответствуют 


друг другу. 


    


    Наконец, когда база геоданных будет исправлена, необходимо избежать 


появления новых ошибок. Таким образом, должны быть созданы процедуры, 


которые могли бы автоматически проверить каждый новый объект, 


вставленный в базу геоданных. 


    


    Обычно ограничения целостности определяют правила пополнения базы 


геоданных. Правила применяются к классу географических объектов и 


заставляют каждый объект этого класса (дорога, например) соответствовать 


этим правилам. Это означает, что объекты реального мира имеют 


предопределенные характеристики. 


    


    Тем не менее, такая гипотеза на практике слишком строгая. Исключения 


всегда имеют место в наборах геоданных. Например, большинство дорог не 


пересекает здания, но некоторые из них проходят внутри здания или под ним. 


Это не означает, что мы не можем использовать такое ограничение, но мы 


должны обеспечить способы обработки исключений. 


    


    Топологические ограничения целостности, которые были описаны выше, 


позволяют обнаруживать топологические ситуации, чтобы исправить их или 


только сообщать о них. Может быть, конечному пользователю необязательно 


исправлять каждое нарушение ограничения, обнаруженное в базе геоданных. 


Он может просто хранить их в базе исключений. 


    


    Существует два способа обработки исключений. Первый должен 


определить каждое ограничение с атрибутом <<{нет исключений}>>  и 


сообщать о каждом случае, как только такое исключение будет обнаружено. 


База исключений тогда пополняется постепенно присоединением списка 


случаев исключений к каждому ограничению. В~этом случае исключение~--- 


это выражение, содержащее пару объектов. Например:


    \smallskip


    


\noindent


Изначально 





\noindent


ОГРАНИЧЕНИЕ~1\;=\;(Дорога, Пересечение, Здание, Запрещенное) 








\noindent


нет исключений. 





\smallskip





\noindent


После процесса обнаружения





\noindent


ОГРАНИЧЕНИЕ~1\;=\;(Дорога, Пересечение, Здание, Запрещенное)  





\noindent


с исключением (Искл24, Сущность47); (Искл53, Сущность47). 


     


\smallskip





    Второй способ заключается в том, чтобы позволить конечному 


пользователю создавать свой список исключений заранее. В~этом случае пары 


объектов не проверяются в течение процесса обнаружения. На иерархической 


структуре определения геообъектов мы можем определить те ограничения 


целостности, которые унаследованы через иерархию и позволяют 


определять исключения на подклассах. Например:


    


    \noindent


Туннель~--- подкласс дороги. 





\noindent


ОГРАНИЧЕНИЕ~2\;=\;(Дорога, Пересечение, Здание, Запрещенное)  





\noindent


с исключением (туннель, здание). 





\section{Заключение}


    


    В работе представлена структура для согласованности и коррекции 


наборов геоданных с учетом трех видов ошибок (структурные, геометрические 


и семантические. Проведение коррекций предполагает выполнение трех этапов: 


    \begin{enumerate}[1.]


    \item \textit{Подготовка к проверке.} В~соответствии с содержанием базы 


геоданных должен быть выбран набор свойств и определен набор 


топологических ограничений целостности.


    \item \textit{Свойства и проверка ограничений.} Свойства проверяются 


посредством вычислительных геометрических алгоритмов, а ограничения  


проверяются предикатами первого порядка. Этот этап выявляет ряд ошибок. 


    \item \textit{Исправление ошибок.} Сценарии коррекции создаются для 


каждого типа ошибки. Семантические (исходящие из ограничений 


топологической целостности) и геометрические ошибки обеспечиваются 


набором соответствующих инструментальных средств. 


    \end{enumerate}


    


    Целесообразно объединять различные подходы к выполнению коррекций, 


чтобы обеспечить различные уровни функционирования в зависимости от 


профиля пользователя. 


    


    Как было описано выше, не следует ограничиваться конкретной моделью 


геоданных (за исключением примеров структурных ошибок). Должна быть 


сформирована открытая система, в которой могут быть учтены новые виды 


ошибок, так как мы не можем убедиться в полноте перечисленных ошибок, 


которые будут найдены в пространственной базе геоданных~[13]. Кроме того, 


возникают новые потребности, с которыми связаны новые ограничения в 


пространственных моделях геоданных (и новые виды ошибок). 


    


    Чтобы применять структуру к конкретной ГИС~[14], необходимо 


произвести следующую адаптацию: 


    \begin{itemize} %[1)]


\item проверка структурной согласованности и процессы коррекции должны 


быть перепроектированы согласно структурам геоданных конкретной ГИС; 


\item для каждой ГИС должен быть установлен свой список свойств; 


\item для каждой ГИС должны быть определены ограничения топологической 


целостности. 


\end{itemize}








{\small\frenchspacing


{%\baselineskip=10.8pt


\addcontentsline{toc}{section}{Литература}


\begin{thebibliography}{99}


\bibitem{1-2-8}


\Au{Van Oort P.}


ISO standards for geographic information. Book review~/\!\!/ Int. J. Appl. Earth 


Observation Geoinformation, 2004. Vol.~6, No.\,2. P.~159--160.





\bibitem{2-2-8}


\Au{Kainz W.}


Logical consistency~/\!\!/ Elements of spatial data quality.~---  Elsevier Science, 


1995. P.~109--138. 





\bibitem{3-2-8}


\Au{Salge P.}


Semantic accuracy~/\!\!/ Elements of spatial data quality.~---  Elsevier Science, 1995. 


P.~139--152. 





\bibitem{4-2-8}


\Au{Billen R., Van de Weghe N.}


Qualitative spatial reasoning~/\!\!/ Int. Encyclopedia Human Geography, 


2009. Vol.~5. P.~12--18.





\bibitem{5-2-8}


\Au{Egenhofer M.\,J.}


Reasoning about binary topological relations~/\!\!/ 2nd Symposium on Large Spatial 


Databases, SSD'91, Proceedings. Lecture Notes in Computer Science, 


1991.Vol.~525. P.~143--159.





\bibitem{6-2-8}


\Au{Egenhofer M.\,J., Herring J.\,R.}


A~mathematical framework for the definition of topological relationships~/\!\!/ 4th 


Symposium (International) on Spatial Data Handling, SDH-90, Proceedings.~--- 


Zurich, 1990. P.~803--813. 





\bibitem{7-2-8}


\Au{Clementini E., Sharma J., Egenhofer~M.\,J.}


Modelling topological spatial relations: Strategies for query 


processing~/\!\!/ Computers Graphics, 1994. Vol.~18. No.\,6.  P.~815--822.





\bibitem{8-2-8}


\Au{Shihong Du, Qimin Qin, Qiao Wang, Haijian Ma}.


Evaluating structural and topological consistency of complex regions with broad 


boundaries in multi-resolution spatial databases~/\!\!/ Information Sci., 2008. 


Vol.~178. No.\,1. P.~52--68.





\bibitem{9-2-8}


\Au{Liu K., Shi W.}


Extended model of topological relations between spatial objects in geographic 


information systems~/\!\!/ Int. J.~Appl.\ Earth Observation Geoinformation, 


2007. Vol.~9. No.\,3. P.~264--275.





\bibitem{10-2-8}


\Au{Clementini E.}


A~model for uncertain lines~/\!\!/ J.~Visual Languages Computing, 2005. 


Vol.~16. No.\,4. P.~271--288.





\bibitem{11-2-8}


\Au{Liu Yu, Guo Q., Kelly M.}


A~framework of region-based spatial relations for non-overlapping features and its 


application in object based image analysis~/\!\!/ ISPRS J.~Photogrammetry Remote 


Sensing, 2008. Vol.~63. No.\,4. P.~461--475.





\bibitem{12-2-8}


\Au{Schockaert S., Smart P.\,D., Twaroch~F.\,A.}


Generating approximate region boundaries from heterogeneous spatial information: 


An evolutionary approach~/\!\!/ Information Sci., 2011. Vol.~181. No.\,2. 


P.~257--283.





\bibitem{13-2-8}


\Au{Дулин С.\,К., Дулина Н.\,Г.}


О~проблеме согласованности базы геоданных.~--- М.: ВЦ РАН, 2007. 21~с.





 \label{end\stat}





\bibitem{14-2-8}


\Au{Розенберг И.\,Н., Дулин С.\,К.}


Геоинформационный портал отрасли. Гарантировать достоверность 


данных~/\!\!/ Железнодорожный транспорт, 2010. №\,2.  С.~12--17.


\end{thebibliography}


}


}    


